\RequirePackage{etex}
\documentclass[11pt, onecolumn, final]{IEEEtran}

%%%%%%%%%%%%%%%%%%%%%%%%%%%%%%%%%%%%%%%%%%%%%%%%%%%%%%%%%%%%%%%%%%%%

\def\maintext{tips} %input your manuscript file name here

\def\docprefix{R} %prefix to use for additional references and other local features

%%%%%%%%%%%%%%%%%%%%%%%%%%%%%%%%%%%%%%%%%%%%%%%%%%%%%%%%%%%%%%%%%%%%

\usepackage{lmodern}
\usepackage{anyfontsize}
\usepackage{expl3}
\usepackage{silence}
\WarningFilter{caption}{Unsupported document class}
\WarningFilter{caption}{Unknown document class}
\WarningFilter{latex}{No \author given}


\usepackage[labeled,resetlabels]{multibib}
% Used by additional references. If you do not have any local cites, comment this out.
% Cite additional references using \citeA. 
\newcites{A}{Additional References}

\usepackage[response]{signalpreamble} 


\myexternaldocument{\maintext}

%%%%%%%%%%%%%%%%%%%%%%%%%%%%%%%%%%%%%%%%%%%%%%%


%%%%%%%%%%%%%%%%%%%%%%%%%%%%%%%%%%%%%%%%%%%%%%%
% Start of response to reviewers
%%%%%%%%%%%%%%%%%%%%%%%%%%%%%%%%%%%%%%%%%%%%%%%


\title{Response to Reviewers' Comments}


\begin{document}
\maketitle

% This copies the new commands and acronyms from the main paper.
\ExecuteMetaData[\maintext]{tag:commands-main}


We sincerely thank the Associate Editor and all reviewers for their helpful comments and suggestions. We have revised the manuscript and included here a point-by-point response. The main changes are the following:
\begin{enumerate}
  \item ...
  \item ...
\end{enumerate}
\medskip

Major changes are highlighted in \blue{blue} in our revised manuscript. The references like equation numbers refer to those in the manuscript, except for those prefixed by ``\docprefix''. 

%%%%%%%%%%%%%%%%%%%%%%%%%%%%%%%%%%%%%%%%%%%%

% Initial a new point-by-point response to a new reviewer
\PbPResponseHeader

\noindent\textbf{Recommendation:} Major Revision

\begin{comment}
A lot of comments.
\end{comment}

\begin{response}
  Thank you for your comment. Please always use the formal ``thank you'' instead of ``thanks''. Please always start your response politely and in a conciliatory tone even if you are frustrated with the comment. We can be firm, friendly and respectful. After all, the reviewer has sacrificed his or her precious time to read our manuscript and provide comments, which he or she is not obligated to do.

  Using the xr package, I can reference \cref{sect:General}, \cref{fig:network} and \cref{eq:X} of the main paper tips.tex. To do this, you need to build response.tex in the same folder containing tips.aux. 

  Next, I cite \cite{LenTayQui:J16,Str:93}, which are in the main paper. The references displayed are from tips.aux and tips.bbl thanks to xr. Finally, the catchfilebetweentags package allows me to quote the following from \cref{sect:Response} in tips.tex:
  \revision{tag:catch}

  Note that the tag name cannot include underscores. Another quote from tips.tex:
  \revision{tag:equation}

  The equation numbers shown here are the same as those in the main paper, which is great!
\end{response}

\begin{comment}
Do something.
\end{comment}
\begin{response}
  Thank you for your suggestion. I will now cite some additional references that appear only in the bibliography for this response using \citeA{AceDahLobOzd:11,ChePap:95,Cov:69} and \citeA[Thm 1]{JiaTayEld:J24}. Note that we use \verb|\citeA| to cite these references, while \verb|\cite| will cite references from the main paper.
\end{response}


% Initial a new point-by-point response to a new reviewer
\PbPResponseHeader

\noindent\textbf{Recommendation:} Major Revision

\begin{comment}
Why does this hold?
\end{comment}
\begin{response}
  Thank you for your question. Please refer to the following equation:
  \begin{align}
  f(x) &= \sum_{i>0} y_i. \label{eq:ftest}
  \end{align}
  Note the prefix ``\docprefix'' in the equation that is in this response. We can reference \cref{eq:ftest} as usual.

\end{response}

\begin{comment}
A lot of comments.
\end{comment}
\begin{response}
  Thank you for your comment. We have revised the manuscript as follows:
  \revision{tag:equation}
  We remark that it can be shown that
  \begin{align}
  \Lambda(k) &= \int_\Omega \tau(b) \ud{\nu(b)}, \label{eq:Ltest}\\
  \Gamma(k) &= \int_{\calA} f(x) \ud{\mu(x)}. \label{eq:Gtest}
  \end{align}
  The equation numbers like \cref{eq:Lambda_SGLR} shown in the box are the same as those in the main paper, which is great! The command \verb|\revision| takes care of this. But it has limitations. See the following examples. The local equations are \cref{eq:Ltest,eq:Gtest}.
\end{response}


\begin{comment}
Some other comments.
\end{comment}
\begin{response}
Thank you for your comment. Please refer to \cref{fig:tandem}. We provide the following example, reproduced here:   
\revision{tag:ex}[Example>eg:subGSP;figure>fig:fractal]

Note the use of \verb|\xrcounterref| in the macro \verb|\revision| to set the numbering of environments correctly. \textbf{If you know how to make this automatic, let me know.}

The macro \verb|\revision| takes in an optional list of pairs of environment name and label, separated by semicolons ;. Each pair is of the form \verb|environment_name>label|. The environment name can be figure, table, Theorem, Lemma, Example, etc. The label is the label of the environment in the main paper. For example, in the above example, we have \verb|Example>eg:subGSP;figure>fig:fractal|.  
\end{response}

\begin{comment}
Some other comments.
\end{comment}
\begin{response}
We provide the following theorem, reproduced here:

\revision{tag:thm}[Theorem>thm:limit_var;Corollary>cor:limit_var]

Note that the theorem and corollary numbers shown in the box are the same as those in the main paper.
\end{response}

\begin{comment}
Some other comments.
\end{comment}
\begin{response}
Thank you for your comment. Blah blah blah blah blah blah blah blah blah blah blah blah blah blah blah blah blah blah blah blah blah blah blah blah blah. 

In the macro \verb|\revision|, floats like tables and figures have their placement options replaced with ``H''. For example, the following figure and tables from the main text are reproduced here:
\revision{tag:afigure}[figure>fig:tandem]
\revision{tag:atable}[table>tab.Frond_com]

The following is a new \cref{tab:ablation} that is not in the main text, but is included here for demonstration purposes (notice that the table counter value is correctly set to I):

\begin{table}[H]
\centering
\caption{Module ablation.}\label{tab:ablation}
\begin{tabular}{l |c c c}
\toprule
Model &   \multicolumn{3}{c}{KITTI360-AG (Sat.)} \\
   &  R@1 & R@5  & R@10 \\
\midrule
full & \first{32.0} & \first{47.6} & \first{54.9} \\
full w/o stage 1 & 28.5 & 43.9 & 51.0 \\
full w/o stage 2 & 30.6 & 45.3 & 53.0 \\
full w/o $\ell_{\mathrm{MVMM}}$ & 30.9  & 45.2 & 51.7 \\
\bottomrule
\end{tabular}
\end{table}

\end{response}


\begin{comment}
Some other comments.
\end{comment}
\begin{response}
See the following example where the tag is nested inside another tag:
\revision{tag:kernelGSP}

\end{response}




%%%%%%%%%%%%%%%%%%%%%%%%%%%%%%%%%%%%%%%%%%%%%%
% Used by \citeA, if you do not have any local cites, comment this out.
\clearpage
\bibliographystyleA{IEEEtran}
\bibliographyA{IEEEabrv,StringDefinitions,refs}


\end{document}
